% 02-fundamentacao.tex -- incluído por main-pt.tex
% =====================================================================
\section{Fundamentação Teórica}
\label{sec:background}
% =====================================================================

\subsection{Mercados de Carbono, REDD+ e o SBCE}

Os mercados de carbono permitem que as emissões evitadas ou removidas por uma atividade sejam quantificadas, certificadas e negociadas como créditos, em que cada crédito corresponde tipicamente a uma tonelada de dióxido de carbono equivalente \citep{Aguiar2018}. Os projetos REDD+ geram créditos ao reduzir o desmatamento e a degradação florestal em relação a um cenário de referência, o que os torna particularmente relevantes para o Brasil \citep{Aguiar2018,Neves2025}. A credibilidade desses projetos depende do processo de mensuração, relato e verificação (MRV), historicamente caracterizado por dados opacos, altos custos de transação, procedimentos manuais e dependência de intermediários, o que torna o mercado vulnerável a fraudes e à dupla contagem \citep{Kotsialou2022,Howson2019geoforum,Caburlao2026fsi}.

A Lei nº 15.042/2024 estabeleceu o SBCE, no qual operadores com emissões anuais acima de $10.000$~tCO$_2$e devem apresentar plano de monitoramento e relatos de emissões, e aqueles acima de $25.000$~tCO$_2$e devem, adicionalmente, conciliar periodicamente suas obrigações, mantendo ativos do SBCE em quantidade equivalente às suas emissões líquidas \citep{Lei15042}. O sistema prevê um registro central gerido por um órgão gestor, a integração de créditos verificados por metodologias registradas e programas jurisdicionais de REDD+ de abrangência nacional ou estadual. Do ponto de vista computacional, dois aspectos dessa estrutura são determinantes: o conjunto de entidades que podem validar registros é institucionalmente delimitado, e as operações tendem a se concentrar ao redor dos prazos de conciliação.

\subsection{Tecnologia Blockchain}

Uma blockchain é um registro distribuído no qual as transações são agrupadas em blocos encadeados criptograficamente e replicados entre os participantes de uma rede ponto a ponto \citep{Nakamoto2008}. Suas propriedades de descentralização, imutabilidade e transparência a tornam adequada para registrar a emissão, a transferência e a aposentadoria de créditos de carbono, bem como para a sua tokenização \citep{Vilkov2023,Kotsialou2021}. Do ponto de vista da segurança, as garantias de um registro baseado em blockchain são comumente analisadas em termos da tríade confidencialidade, integridade e disponibilidade \citep{Aminzade2018,Tchernykh2019,Rashidibajgan2024}, juntamente com riscos específicos da camada de aplicação, como vulnerabilidades de contratos inteligentes \citep{GhiyamiPour2025}, a confiabilidade dos oráculos que trazem dados de MRV para o registro \citep{Rahman2023} e a governança na própria cadeia \citep{Esposito2025}. A seleção de uma tecnologia de registro distribuído envolve, portanto, compromissos entre desempenho, segurança e descentralização \citep{Werth2023,Rodrigues2024}.

\subsection{Protocolos de Consenso e Complexidade de Comunicação}

Os protocolos de consenso podem ser classificados em famílias baseadas em prova e baseadas em votação \citep{Xiao2020,Bamakan2020}. No consenso de Nakamoto, baseado em prova de trabalho (PoW), o direito de anexar um bloco é obtido pelo dispêndio de um recurso computacional, e a finalidade é probabilística \citep{Nakamoto2008}. Na prova de participação (PoS), esse direito é proporcional à participação dos validadores, e a finalidade pode ser provida por uma camada adicional como o Casper FFG, que finaliza blocos em pontos de verificação periódicos e tolera menos de um terço de participação maliciosa \citep{Buterin2017}. A prova de autoridade com participação (PoSA), modelada como Parlia no JABS, combina um conjunto restrito de validadores autorizados com seleção baseada em participação \citep{Caburlao2026clei}.

Os protocolos baseados em votação derivam da literatura de replicação de máquinas de estado tolerante a falhas bizantinas (BFT) \citep{Vukolic2015}. O pBFT \citep{Castro1999} garante finalidade determinística desde que no máximo $f$ de $n \geq 3f+1$ validadores sejam faltosos. Em seu caso normal, sem falhas, cada bloco requer uma fase de pré-preparação, em que o líder envia $n-1$ mensagens, e duas fases de difusão de todos para todos (preparação e confirmação), o que resulta em um número de mensagens por bloco da ordem de
\begin{equation}
M_{\mathrm{pBFT}}(n) \approx 2n^{2},
\label{eq:mpbft}
\end{equation}
enquanto cada nó envia e recebe $O(n)$ mensagens. O Tendermint \citep{Buchman2018} adota um proponente rotativo e a difusão de votos por \textit{gossip}, preservando duas rodadas de votação de todos para todos. O HotStuff \citep{Yin2019}, por sua vez, concentra os votos no líder e utiliza assinaturas de limiar, o que reduz o número total de mensagens no caso normal para $O(n)$; entretanto, o líder continua a receber e a enviar $O(n)$ mensagens por fase. Essa distinção é central para este trabalho: um protocolo pode ter custo total linear e, ainda assim, apresentar um nó cujo enlace se torna a restrição vinculante.

O caso normal não esgota a questão. \citet{Dolev1985} demonstraram que protocolos determinísticos de acordo bizantino requerem $\Omega(n + f^{2})$ mensagens nas execuções com até $f$ processos faltosos, o que, para $f$ proporcional a $n$, corresponde a $\Omega(n^{2})$. Assim, mesmo um protocolo linear no caso normal, como o HotStuff, incorre em custo quadrático quando sucessivos líderes faltosos forçam trocas de visão. Por essa razão, a avaliação de viabilidade deve incluir validadores bizantinos, e não apenas o caso normal. Como alternativa estrutural, arquiteturas hierárquicas executam o consenso em grupos locais e apenas representantes participam do consenso global \citep{Amir2010}.

\subsection{Avaliação de Desempenho de Sistemas Computacionais}

Segundo \citet{Jain1991}, um estudo sistemático de desempenho compreende a definição dos objetivos e das fronteiras do sistema, a seleção de métricas, a caracterização da carga, a seleção de fatores e níveis, o projeto de experimentos e a análise dos dados. A ANOVA permite testar efeitos e interações e alocar a variação de uma resposta entre os fatores, suas interações e o erro experimental \citep{Jain1991,Montgomery2017}. A parcela de variação atribuída a cada fator, contudo, depende da amplitude dos níveis escolhidos, e não deve ser interpretada como importância intrínseca quando os fatores têm faixas arbitrárias. Quando a avaliação é conduzida por simulação, o erro experimental é controlado por replicações independentes e pela remoção do transiente inicial \citep{Law2015}, e modelos de regressão ajustados aos resultados (metamodelos) só devem ser usados para extrapolação após validação em pontos não utilizados no ajuste \citep{Kleijnen2015}. Neste trabalho, essa metodologia sustenta a avaliação de um artefato no paradigma de \textit{design science} \citep{Hevner2004}: um método de dimensionamento de registros de carbono, composto pela seleção de protocolos por requisitos, pelo modelo de demanda, pelas fronteiras de viabilidade e pela regra de recomendação.
